\section{Original Datasets}
\label{sec:original}

Both benchmarks are derived by \citet{zhou2024grasping}. \textbf{DefOn-FewRel} comes from FewRel \citep{han2018fewrel}, whose public portion has 80 relations with 700 crowd-annotated instances each. \textbf{DefOn-WikiZSL} comes from Wiki-ZSL \citep{chen2021zs}, with 113 relations and 93,383 distantly supervised instances. Relations are Wikidata properties; the authors rewrote property definitions into sentences with entity placeholders. They sampled 5 groups of 14 FewRel relations and 3 groups of 15 Wiki-ZSL relations as test sets.

The paper states that the unlabelled corpus for each group is 10,000 instances down-sampled from the 899,996-instance corpus of \citet{gao2020neural}. The released splits differ, as Table~\ref{tab:splits} shows: every released \texttt{distant.json} has 100,000 instances. Section~\ref{sec:replication} returns to this.

\begin{table}[t]
\centering\small
\begin{tabular}{lrrr}
\toprule
Dataset & Rel. & Test (per split) & Pool \\
\midrule
FewRel (5 splits) & 14 & 9,796--9,799 & 100,000 \\
WikiZSL (3 splits) & 15 & 12,122--12,857 & 100,000 \\
\bottomrule
\end{tabular}
\caption{Released splits, counted from the JSON files.}
\label{tab:splits}
\end{table}

Each instance holds tokens plus head and tail mentions with token indices. Each relation has a definition prompt. Real examples from the test files follow.

\begin{quote}\small
\textbf{FewRel, P106 (occupation).} Definition: \snip{<ENT1> was/is the occupation of <ENT0> (a person)}.
Sentence: ``David Malo or Davida Malo ( 1793--1853 ) was a leading Native Hawaiian historian of the Kingdom of Hawaii .''
Head: \emph{david malo}; tail: \emph{historian}.

\textbf{WikiZSL, P101 (field of work).} Definition: \snip{<ENT1> (a field of work, academic area, or specialism) was/is the specialization of <ENT0> (a person or organization)}.
Sentence: ``Otto Toeplitz ( 1 August 1881 15 February 1940 ) was a German Jewish mathematician working in functional analysis .''
Head: \emph{Otto Toeplitz}; tail: \emph{functional analysis}.
\end{quote}

The authors attribute WikiZSL's lower scores to noise from distant supervision. These datasets are suitable because their definitions are precise and they have enough relations to average over, though the FewRel test instances are drawn from a single domain (Wikipedia).
